\thispagestyle{empty}
\section*{Аннотация}

В работе разработана и реализована система <<AI Coauthor>> — комплекс семантического
анализа научного пространства, тематической кластеризации публикаций и рекомендации
потенциальных научных соавторов с объяснениями на основе больших языковых моделей.
Цель работы — построить инструмент, который по тексту научных статей выявляет тематическую
структуру предметной области и обоснованно подсказывает исследователю семантически
совместимых коллег. Эмпирическую базу составляет корпус из 509\,447 публикаций OpenAlex
по подполю Artificial Intelligence и 35\,954 квалифицированных авторов. Статьи представлены
многоязычными эмбеддингами модели multilingual-e5-base (768 измерений); тематическая
структура из 979 тем получена пайплайном в стиле BERTopic (UMAP $\to$ HDBSCAN $\to$
c-TF-IDF $\to$ MMR $\to$ разметка языковой моделью), а 979 тем дополнительно объединены
в 76 метакластеров — семантический слой над темами. HDBSCAN сравнивался с наивными
бейзлайнами K-Means и случайным разбиением на UMAP-10; выбор обоснован по совокупности
независимых метрик качества тем. Рекомендация соавторов реализована в два этапа:
mean-, Transformer- и GraphSAGE-представления автора извлекают близких кандидатов (retrieval),
после чего пятиступенчатый пайплайн на основе языковой модели переранжирует их и формирует
объяснения с конкретными попарными связями статей. Качество финальной выдачи оценено
LLM-судьёй: LLM Has3@10 (score $\geq 2$) = 0,875, LLM Selected3Mean@10 = 2,49 из 3.
Результат интегрирован в
веб-приложение с интерактивной картой научного пространства.

\newpage
\thispagestyle{empty}
\section*{Abstract}

\begin{english}
This work presents AI Coauthor, a system for the semantic analysis of the scientific landscape:
it performs topic clustering of publications and recommends potential research collaborators with
explanations produced by large language models. The goal is a tool that recovers the topical
structure of a research field directly from paper texts and suggests semantically compatible peers
in a justified, interpretable way. The empirical base is a corpus of 509{,}447 OpenAlex publications
in the Artificial Intelligence subfield and 35{,}954 qualified authors. Papers are represented by
multilingual embeddings from the multilingual-e5-base model (768 dimensions); a topical structure
of 979 topics is obtained with a BERTopic-style pipeline (UMAP $\to$ HDBSCAN $\to$ c-TF-IDF
$\to$ MMR $\to$ LLM labeling), and the 979 topics are further grouped into 76 meta-clusters, a
semantic layer over the topics. HDBSCAN is compared against naive K-Means and random baselines on
UMAP-10; the choice is justified through a set of independent topic-quality metrics. Collaborator
recommendation is built in two stages: mean, Transformer, and GraphSAGE author representations retrieve nearby
candidates, after which a five-stage LLM pipeline re-ranks them and generates explanations grounded
in concrete pairwise paper links. The final recommendations are evaluated with an LLM judge:
LLM Has3@10 (score $\geq 2$) = 0.875, and LLM Selected3Mean@10 reaches 2.49 out of 3. The
result is integrated into a web application with an
interactive map of the scientific space.
\end{english}
